%==============================================================================
% house-common.tex -- the preamble shared by documents and slides
%
% Loaded first by house-pdf and house-beamer, ahead of house.tex or
% house-beamer.tex. It holds what a page and a slide treat alike: tables, the
% math sans alphabet, code and printed output.
%==============================================================================

%-------------------------------------------------------------- math alphabets
% Expectation is written \mathsf{P}. By default unicode-math sets \mathsf in the
% text sans family, which is Lato, the heading face; in a document that named no
% sans it was Latin Modern Sans. `sym` takes the letter from the math font's own
% sans-serif alphabet instead, which STIX Two Math draws to match the rest of
% the mathematics. Guarded, because unicode-math is loaded only when a math font
% is named.
\makeatletter
\@ifundefined{unimathsetup}{}{\unimathsetup{mathsf=sym}}
\makeatother

%---------------------------------------------------------------------- tables
% greater_tables 7 draws a table with its `tabular` emitter: one real tabular
% inside a single tikz node, taking its horizontal rules from booktabs and
% drawing only the verticals in tikz. Pandoc's template loads booktabs and
% array only when the document happens to contain a markdown table, which is
% not a thing to depend on. The emitter defines its own \I strut with
% \providecommand, so the preamble no longer needs to.
\usepackage{booktabs}
\usepackage{array}
\usepackage{xcolor}
\usepackage{tikz}

%------------------------------------------------------------------------ code
% \codesize is the one dial for code and printed output. Documents set code at
% \footnotesize; house-beamer.tex redefines the dial for the slide canvas.
% \fvset stores the macro, not its meaning, so the later definition wins.
\providecommand{\codesize}{\footnotesize}

% Smaller code, and long lines break rather than running off the page.
% \fvset sets the defaults every fancyvrb environment inherits, which includes
% pandoc's Highlighting, so Quarto's own options on it survive; redefining the
% environment would discard them. breakanywhere is needed because Python has
% long dotted names with no break points.
\usepackage{fvextra}
\fvset{breaklines=true,breakanywhere=true,breaksymbolleft={},fontsize=\codesize}

% Printed output (a cell's stdout, and the text tables `qd` prints) arrives as
% LaTeX's own verbatim, which \fvset cannot reach, so it is rerouted through
% fancyvrb and takes the size and the line breaking above.
%
% \AtBeginDocument, not here, and that is load bearing. Quarto appends its
% callout machinery after include-in-header, and loading tcolorbox there pulls
% in the classic verbatim environment again, silently undoing a redefinition
% made at this point in the preamble. The symptom is a cell's stderr running
% off the right edge of the page: a warning carrying an absolute path has no
% spaces to break at, so breakanywhere is the only thing holding it on the
% page. Deferring puts this definition last, where it wins.
\AtBeginDocument{\DefineVerbatimEnvironment{verbatim}{Verbatim}{}}
